\chapter{Maxwell's Equations: A Working Review}
\label{ch:maxwell}

% Chapter-local vector macros not in style/simbook.sty (\vs is the time-domain
% state vector; the Poynting vector, wavevector, and z-hat are local here).
\newcommand{\vS}{\mathbf{S}}
\providecommand{\vk}{\mathbf{k}}
\providecommand{\vz}{\mathbf{z}}

\chapterorient{You have met Maxwell's equations before---probably as four lines
on a slide, admired and then set aside. This chapter brings them back as
\emph{working} equations: the ones the simulator actually discretizes. We will
read each of the four laws for its physical content, attach the materials
through the constitutive relations, pin down what happens at a wall or an
interface, account for the energy, and---most important for everything that
follows---split the physics into the three regimes (static, time-harmonic,
time-domain) that become the analyses of Part~II onward. By the end you will
recognize, in raw continuous form, every operator the rest of the book learns
to assemble and solve.}

\section{The four laws, read for meaning}

Electromagnetism is the statement that two vector fields---the electric field
$\vE$ and the magnetic field $\vB$---fill all of space, push on charges, and
generate each other. Maxwell's four equations say exactly \emph{how} the fields
arrange themselves around their sources and how a change in one provokes the
other. In differential (point-by-point) form, in a region of space with charge
density $\rho$ and current density $\vJ$, they are

\begin{align}
  \Div \vD &= \rho, & &\text{(Gauss's law for $\vE$)} \label{eq:gauss-e}\\
  \Div \vB &= 0, & &\text{(Gauss's law for $\vB$)} \label{eq:gauss-b}\\
  \Curl \vE &= -\,\partial_t \vB, & &\text{(Faraday's law)} \label{eq:faraday}\\
  \Curl \vH &= \vJ + \partial_t \vD. & &\text{(Amp\`ere--Maxwell law)} \label{eq:ampere}
\end{align}

Here $\vD$ is the electric displacement and $\vH$ the magnetic field intensity;
they are the ``material-aware'' partners of $\vE$ and $\vB$, related to them by
the constitutive relations of \cref{sec:constitutive}. For the moment read
$\vD$ as ``$\vE$, scaled by how the material responds'' and $\vH$ as
``$\vB$, scaled the same way.'' The operators $\Div$ and $\Curl$ are the
ordinary divergence and curl of vector calculus; if they feel rusty,
\cref{tab:vectorcalc} is a one-line reminder of what each measures.

\begin{table}[t]
\centering
\caption[The two derivative operators]{The two first-order vector-calculus
operators that build every Maxwell equation. \emph{Divergence} measures
net outflow from a point (a source); \emph{curl} measures local circulation
(a swirl). Maxwell's equations are entirely statements about the divergence
and curl of $\vE$ and $\vB$.}
\label{tab:vectorcalc}
\small
\setlength{\tabcolsep}{6pt}
\renewcommand{\arraystretch}{1.25}
\begin{tabular}{@{}lll@{}}
\toprule
Operator & Reads & Measures \\
\midrule
$\Div \vF$ & ``div of $\vF$'' & net flux out of an infinitesimal box (a source/sink) \\
$\Curl \vF$ & ``curl of $\vF$'' & circulation around an infinitesimal loop (a swirl) \\
\bottomrule
\end{tabular}
\end{table}

A field theory is much easier to feel than to memorize. Take the four laws one
at a time, each with the picture it draws (\cref{fig:fourlaws}).

\begin{figure}[t]
\centering
\begin{tikzpicture}[scale=1.0]
% --- Gauss for E: charge with radial field lines ---
\begin{scope}[shift={(0,0)}]
  \fill[simRust] (0,0) circle (2.2pt);
  \node[font=\scriptsize] at (0.0,-0.28) {$+q$};
  \foreach \a in {0,45,...,315}{
    \draw[flow,simBlue,shorten <=3pt]
      (0,0) -- (\a:0.95);}
  \node[font=\scriptsize\itshape,simBlue] at (0.85,0.85) {$\vE$};
  \node[font=\footnotesize] at (0,-1.35) {$\Div\vD=\rho$};
  \node[font=\scriptsize,text=simGray,align=center] at (0,-1.75)
    {charge is\\ a source};
\end{scope}
% --- Gauss for B: dipole loop, no source ---
\begin{scope}[shift={(3.6,0)}]
  \draw[flow,simTeal] (0,0) circle (0.75);
  \draw[flow,simTeal] (0,0) circle (0.42);
  \node[font=\scriptsize\itshape,simTeal] at (0.95,0.0) {$\vB$};
  \fill[simInk] (0,0) circle (1.2pt);
  \node[font=\footnotesize] at (0,-1.35) {$\Div\vB=0$};
  \node[font=\scriptsize,text=simGray,align=center] at (0,-1.75)
    {lines close;\\ no source};
\end{scope}
% --- Faraday: changing B makes circulating E ---
\begin{scope}[shift={(7.2,0)}]
  \draw[-{Stealth[length=2mm]},simRust,thick] (0,-0.55)--(0,0.55)
    node[right,font=\scriptsize]{$\partial_t\vB$};
  \draw[flow,simBlue] (0.85,0) arc (0:300:0.85);
  \node[font=\scriptsize\itshape,simBlue] at (-0.95,0.5) {$\vE$};
  \node[font=\footnotesize] at (0,-1.35) {$\Curl\vE=-\partial_t\vB$};
  \node[font=\scriptsize,text=simGray,align=center] at (0,-1.75)
    {changing $\vB$\\ swirls $\vE$};
\end{scope}
% --- Ampere-Maxwell: current makes circulating H ---
\begin{scope}[shift={(10.8,0)}]
  \draw[-{Stealth[length=2mm]},simGreen,thick] (0,-0.55)--(0,0.55)
    node[right,font=\scriptsize]{$\vJ$};
  \draw[flow,simViolet] (0.85,0) arc (0:300:0.85);
  \node[font=\scriptsize\itshape,simViolet] at (-0.95,0.5) {$\vH$};
  \node[font=\footnotesize] at (0,-1.35) {$\Curl\vH=\vJ+\partial_t\vD$};
  \node[font=\scriptsize,text=simGray,align=center] at (0,-1.75)
    {current swirls\\ $\vH$};
\end{scope}
\end{tikzpicture}
\caption[The four laws as pictures]{Each Maxwell equation as a field-line
sketch. \emph{Gauss for $\vE$}: field lines begin and end on charge.
\emph{Gauss for $\vB$}: field lines form closed loops---there is no magnetic
charge. \emph{Faraday}: a time-varying magnetic field wraps an electric field
around it. \emph{Amp\`ere--Maxwell}: a current (or a time-varying $\vD$) wraps a
magnetic field around it. The last two are mirror images with a crucial sign
difference.}
\label{fig:fourlaws}
\end{figure}

\subsection{Gauss's law: charge is the source of \texorpdfstring{$\vE$}{E}}

Equation~\eqref{eq:gauss-e}, $\Div\vD=\rho$, says that electric field lines
\emph{begin and end on charge}. The divergence of $\vD$ at a point is the net
number of field lines streaming out of an infinitesimal box around that
point; the equation pins that outflow to the charge density enclosed. Where
there is positive charge, lines emanate; where negative, they terminate; in
empty space the divergence is zero and lines merely pass through. Integrate
over a closed surface and you recover the familiar integral form,
$\oint \vD\cdot\vn\,\mathrm{d}S = Q_{\text{enc}}$: the flux through any closed
surface equals the charge inside. This is the law that makes a capacitor a
capacitor, and it is the entire content of the electrostatic analysis.

\subsection{Gauss's law for \texorpdfstring{$\vB$}{B}: no magnetic charge}

Equation~\eqref{eq:gauss-b}, $\Div\vB=0$, is the same statement with the
right-hand side set permanently to zero: there is no magnetic charge. Magnetic
field lines never begin or end; they always close on themselves
(\cref{fig:fourlaws}, second panel). This innocent-looking zero is
load-bearing for the discretization. It forces $\vB$ to be \emph{divergence
free}, and a finite-element method that fails to honour that property produces
spurious, unphysical modes. We will spend real effort in \cref{ch:bc} building
function spaces that respect $\Div\vB=0$ exactly rather than approximately.

\subsection{Faraday's law: a changing \texorpdfstring{$\vB$}{B} drives \texorpdfstring{$\vE$}{E}}

Equation~\eqref{eq:faraday}, $\Curl\vE=-\partial_t\vB$, is the first of the two
\emph{coupling} laws---the ones that knit $\vE$ and $\vB$ together into a single
electromagnetic field. It says a magnetic field that changes in time wraps an
electric field around itself. The curl of $\vE$---its local circulation---is
fed by the rate of change of $\vB$ threading through. This is the law of the
electric generator and the transformer: move a magnet, and the changing flux
induces a circulating $\vE$ that drives current. The minus sign is Lenz's law
in disguise; it ensures the induced effect \emph{opposes} the change that
caused it, which is why energy is conserved rather than spontaneously created.

\subsection{Amp\`ere--Maxwell: current (and changing \texorpdfstring{$\vD$}{D}) drives \texorpdfstring{$\vH$}{H}}

Equation~\eqref{eq:ampere}, $\Curl\vH=\vJ+\partial_t\vD$, is Faraday's mirror.
A current $\vJ$ wraps a magnetic field around itself---the law of the
electromagnet. Maxwell's own contribution was the second term, the
\emph{displacement current} $\partial_t\vD$: even with no moving charge, a
\emph{changing} electric field sources $\vH$ exactly as a real current would.
That term is what lets the two coupling laws close a loop---$\vE$ feeds $\vB$
through Faraday, $\vB$ feeds $\vE$ through Amp\`ere--Maxwell---and that loop,
propagating through space, \emph{is} light. Without the displacement current,
electromagnetic waves do not exist.

We can make ``that loop is light'' precise in three lines, and it is worth
doing once. Take the curl of Faraday's law~\eqref{eq:faraday} and feed in
Amp\`ere--Maxwell~\eqref{eq:ampere} (in vacuum, $\vJ=\bm0$,
$\vD=\varepsilon_0\vE$, $\vB=\mu_0\vH$):
\[
  \Curl(\Curl\vE)
    = -\,\partial_t(\Curl\vB)
    = -\,\mu_0\varepsilon_0\,\partial_t^2\vE.
\]
The vector identity $\Curl\Curl\vE=\Grad(\Div\vE)-\Lap\vE$, together with
$\Div\vE=0$ in charge-free vacuum, collapses the left side to $-\Lap\vE$, and we
are left with the \emph{wave equation}
\begin{equation}
  \Lap\vE - \mu_0\varepsilon_0\,\partial_t^2\vE = \bm0. \label{eq:waveeq}
\end{equation}
This is the equation of a disturbance propagating at speed
$c=1/\sqrt{\mu_0\varepsilon_0}$---and the only thing it took to produce it was
the displacement current, which supplied the second time derivative. Maxwell's
quiet extra term turned a pair of static field laws into a dynamical theory of
waves. The same operator $\Curl\Curl-\,(\text{constant})\,\partial_t^2$ is, in
the frequency domain, the curl--curl operator $\mat K-\omega^2\mat M$ of
Worked example~\ref{wex:phasor}; we will meet it again and again.

\begin{keyidea}
The four laws split cleanly into two pairs. The two \emph{Gauss} laws are
constraints on the field at one instant---they say where lines can begin and
end. The two \emph{curl} laws are the dynamics---they say how $\vE$ and $\vB$
generate each other in time. Static analyses keep only the constraints; wave
and resonance analyses turn on the coupling. Every analysis in this book is a
choice of which terms of \eqref{eq:gauss-e}--\eqref{eq:ampere} survive.
\end{keyidea}

\subsection{The continuity equation: charge is conserved}

The four laws are not independent; they already contain the conservation of
charge. Take the divergence of the Amp\`ere--Maxwell law~\eqref{eq:ampere}. The
divergence of any curl is identically zero, so the left side vanishes:
\[
  0 = \Div(\Curl\vH) = \Div\vJ + \partial_t(\Div\vD).
\]
Substitute Gauss's law $\Div\vD=\rho$ into the last term and you obtain the
\emph{continuity equation}
\begin{equation}
  \Div\vJ + \partial_t\rho = 0. \label{eq:continuity}
\end{equation}
In words: the net current flowing out of any region equals the rate at which
the charge inside it drops. Charge is never created or destroyed---it only
flows. This is not an extra assumption; it is a \emph{consequence} of Maxwell's
equations, and the displacement current is precisely the term that makes it
come out right. (Had Maxwell omitted $\partial_t\vD$, the divergence of
\eqref{eq:ampere} would have forced $\Div\vJ=0$, false whenever charge piles
up.) We will meet the identity $\Div(\Curl(\cdot))=0$ again as one half of the
\emph{de Rham complex} in \cref{ch:derham}; it is the same algebraic fact wearing
different clothes.

\begin{pitfall}
It is tempting to treat $\Div\vJ+\partial_t\rho=0$ as a fifth law to be imposed
separately. It is not---it follows from the other four. In the discrete setting
the analogue is subtle: a finite-element scheme that discretizes the curl laws
\emph{without} a compatible discretization of the divergence will generally
\emph{violate} discrete continuity, leaking charge. Honouring these identities
\emph{discretely}, not just continuously, is the entire reason for the
structured function spaces of \cref{ch:functionspaces}.
\end{pitfall}

\subsection{The integral forms, and why we keep the differential ones}
\label{sec:integral}

Each differential law has an \emph{integral} twin, obtained by integrating over
a region and applying the two great theorems of vector calculus---the
divergence theorem (for the Gauss laws) and Stokes' theorem (for the curl
laws). The divergence theorem turns a volume integral of a divergence into a
flux through the bounding surface; Stokes' theorem turns a surface integral of a
curl into a circulation around the bounding loop (\cref{fig:theorems}). Applied
to the four laws they give

\begin{align}
  \oint_{\partial \mathcal V} \vD\cdot\vn\,\mathrm{d}S &= \int_{\mathcal V}\rho\dV
    = Q_{\text{enc}}, & &\text{(Gauss, integral)}\\
  \oint_{\partial \mathcal V} \vB\cdot\vn\,\mathrm{d}S &= 0, & &\text{(no monopole)}\\
  \oint_{\partial \mathcal S} \vE\cdot\mathrm{d}\bm\ell
    &= -\frac{\mathrm d}{\mathrm dt}\int_{\mathcal S}\vB\cdot\vn\,\mathrm{d}S,
    & &\text{(Faraday: emf $=-\dot\Phi$)}\\
  \oint_{\partial \mathcal S} \vH\cdot\mathrm{d}\bm\ell
    &= \int_{\mathcal S}\Big(\vJ+\partial_t\vD\Big)\cdot\vn\,\mathrm{d}S.
    & &\text{(Amp\`ere--Maxwell)}
\end{align}

These are the forms a physicist reaches for first: ``flux through a closed
surface equals enclosed charge,'' ``the emf around a loop equals minus the rate
of change of flux through it.'' They are equivalent to the differential laws
wherever the fields are smooth---the differential and integral forms are two
readings of the same content, related by the integral theorems. So why does a
simulator prefer the \emph{differential} forms? Because the finite-element
method discretizes a \emph{local} statement: it asks the equation to hold,
weakly, on each little element, and a differential law is already local. The
integral forms remain the right tool for \emph{deriving} the interface
conditions of \cref{sec:bc} (shrink a loop or a pillbox across the interface)
and for \emph{checking} a computed field (does the flux through this surface
equal the charge we put inside?), and we will use them for exactly those two
purposes.

\begin{figure}[t]
\centering
\begin{tikzpicture}[scale=1.0]
% --- divergence theorem: volume with outward flux ---
\begin{scope}[shift={(0,0)}]
  \draw[simInk,thick] (0,0) circle (1.0);
  \fill[simBgBlue,opacity=0.5] (0,0) circle (1.0);
  \node[font=\scriptsize] at (0,0) {$\mathcal V$};
  \foreach \a in {20,90,160,230,300}{
    \draw[-{Stealth[length=1.8mm]},simBlue,thick] (\a:1.0)--(\a:1.45);}
  \node[font=\scriptsize,align=center] at (0,-1.55)
    {$\displaystyle\int_{\mathcal V}\!\Div\vF\,\mathrm{d}V
      =\!\oint_{\partial\mathcal V}\!\vF\cdot\vn\,\mathrm{d}S$};
  \node[font=\footnotesize,text=simGray] at (0,1.7) {divergence theorem};
\end{scope}
% --- Stokes theorem: surface with circulation on boundary ---
\begin{scope}[shift={(5.4,0)}]
  \fill[simBgGold,opacity=0.6] (0,0) ellipse (1.05 and 0.7);
  \draw[flow,simRust] (1.05,0) arc (0:355:1.05 and 0.7);
  \node[font=\scriptsize] at (0,0) {$\mathcal S$};
  \draw[-{Stealth[length=1.8mm]},simGreen,thick] (0,0)--(0,0.55)
    node[right,font=\scriptsize]{$\Curl\vF$};
  \node[font=\scriptsize,align=center] at (0,-1.55)
    {$\displaystyle\int_{\mathcal S}\!\Curl\vF\cdot\vn\,\mathrm{d}S
      =\!\oint_{\partial\mathcal S}\!\vF\cdot\mathrm d\bm\ell$};
  \node[font=\footnotesize,text=simGray] at (0,1.7) {Stokes' theorem};
\end{scope}
\end{tikzpicture}
\caption[The two integral theorems]{The divergence theorem (left) converts the
volume integral of a divergence into the outward flux through the bounding
surface; it turns the differential Gauss laws into their integral forms.
Stokes' theorem (right) converts the surface integral of a curl into the
circulation around the bounding loop; it turns Faraday's and Amp\`ere's laws
into their integral forms. The same two theorems, applied to a test function
instead of the field itself, produce the \emph{weak form} of \cref{ch:weak}.}
\label{fig:theorems}
\end{figure}

These two theorems are not a one-time convenience. The integration by parts that
produces the finite-element \emph{weak form} in \cref{ch:weak} is the
divergence theorem applied to a product of the field and a test function---the
very same identity, used to move a derivative off the field and onto the test
function. The integral forms are the first appearance of the machinery the whole
discretization rests on.

\section{Materials: the constitutive relations}
\label{sec:constitutive}

Maxwell's four laws contain four fields---$\vE,\vD,\vB,\vH$---but only two are
truly independent. The other two are tied to them by the \emph{material}: the
\emph{constitutive relations} encode how a given substance responds to a field.
For a linear, isotropic medium they are three scalar multiplications,

\begin{align}
  \vD &= \varepsilon\,\vE, & &\text{(permittivity: electric response)} \label{eq:const-eps}\\
  \vB &= \mu\,\vH, & &\text{(permeability: magnetic response)} \label{eq:const-mu}\\
  \vJ &= \sigma\,\vE. & &\text{(conductivity: Ohm's law)} \label{eq:const-sigma}
\end{align}

The three coefficients are the only place the \emph{stuff} of the device enters
the physics. Everything else is geometry and excitation.

\begin{itemize}
  \item The \textbf{permittivity} $\varepsilon$ measures how much a material
  polarizes---how its bound charges shift---in response to $\vE$, and so how
  much electric flux $\vD$ a given field produces. In vacuum
  $\varepsilon=\varepsilon_0\approx\SI{8.854e-12}{\farad\per\meter}$. For a
  dielectric one writes $\varepsilon=\varepsilon_r\varepsilon_0$ with the
  \emph{relative permittivity} $\varepsilon_r$ (the ``dielectric constant'')
  ranging from $\approx 1$ (air) through $\approx 3.9$ (silicon dioxide) and
  $\approx 11.7$ (silicon) to $\approx 80$ (water).

  \item The \textbf{permeability} $\mu$ plays the same role magnetically. In
  vacuum $\mu=\mu_0=4\pi\times10^{-7}\,\si{\henry\per\meter}$. Most materials of
  interest in microwave engineering are non-magnetic, $\mu_r\approx 1$;
  ferromagnets ($\mu_r$ in the hundreds or thousands) are the exception. In the
  magnetostatic analysis it is the \emph{reluctivity} $\nu=1/\mu$ that appears
  in the operator, exactly as $\varepsilon$ appears in the electrostatic one.

  \item The \textbf{conductivity} $\sigma$ measures how freely charge flows:
  Ohm's law~\eqref{eq:const-sigma} says the current is proportional to the
  field. A perfect insulator has $\sigma=0$; a perfect conductor has
  $\sigma\to\infty$; copper sits at $\approx\SI{6e7}{\siemens\per\meter}$. The
  conductivity is what makes a material \emph{lossy}---it is the term that
  drains electromagnetic energy into heat, and (as we will see in
  \cref{sec:regimes}) it is what gives a resonant mode a finite quality factor
  $Q$ rather than ringing forever.
\end{itemize}

\begin{table}[t]
\centering
\caption[Typical material parameters]{Representative values of the three
constitutive coefficients. Relative permittivity $\varepsilon_r$ and relative
permeability $\mu_r$ are dimensionless ratios to the vacuum values
$\varepsilon_0\approx\SI{8.854e-12}{\farad\per\meter}$ and
$\mu_0=4\pi\times10^{-7}\,\si{\henry\per\meter}$; conductivity $\sigma$ is in
$\si{\siemens\per\meter}$. The enormous range of $\sigma$---from an insulator
to a metal---is what separates a dielectric from a conductor, and what makes a
boundary a PEC wall in the limit.}
\label{tab:materials}
\small
\setlength{\tabcolsep}{7pt}
\renewcommand{\arraystretch}{1.2}
\begin{tabular}{@{}lccc@{}}
\toprule
Material & $\varepsilon_r$ & $\mu_r$ & $\sigma\;(\si{\siemens\per\meter})$ \\
\midrule
Vacuum             & $1$     & $1$    & $0$ \\
Air                & $1.0006$& $1$    & $\approx 0$ \\
Silicon dioxide (\(\mathrm{SiO_2}\)) & $3.9$ & $1$ & $\sim 10^{-12}$ \\
Silicon            & $11.7$  & $1$    & varies (doping) \\
Water (\SI{20}{\celsius}) & $\approx 80$ & $1$ & $\sim 5\times10^{-4}$ \\
Copper             & ---     & $1$    & $5.96\times10^{7}$ \\
Iron (soft)        & ---     & $\sim 5000$ & $1.0\times10^{7}$ \\
\bottomrule
\end{tabular}
\end{table}

\begin{notationbox}
Three symbols, three responses: $\varepsilon$ couples $\vE\!\to\!\vD$ (electric),
$\mu$ couples $\vH\!\to\!\vB$ (magnetic), $\sigma$ couples $\vE\!\to\!\vJ$
(conduction). Their reciprocals appear in the operators: reluctivity
$\nu=1/\mu$ weights the curl--curl term, and $1/\varepsilon$ never appears
alone---$\varepsilon$ itself is the weight in the electrostatic stiffness. Keep
the pairing $(\vD,\vE,\varepsilon)$ and $(\vB,\vH,\mu)$ straight: $\vD$ and
$\vB$ are the ``flux'' fields (the ones whose divergence the Gauss laws
constrain); $\vE$ and $\vH$ are the ``intensity'' fields (the ones whose curl
the dynamic laws constrain).
\end{notationbox}

\subsection{Isotropic, anisotropic, and the material tensor}

Equations~\eqref{eq:const-eps}--\eqref{eq:const-sigma} assume the material
responds the \emph{same in every direction}---it is \emph{isotropic}, and the
coefficients are scalars. Real materials need not be so accommodating. A
crystal, a layered substrate, or a magnetized ferrite responds differently
along different axes, and then $\varepsilon$ becomes a $3\times3$ symmetric
\emph{tensor} $\bm{\varepsilon}$:
\[
  \vD = \bm{\varepsilon}\,\vE,
  \qquad
  \bm{\varepsilon} =
  \begin{pmatrix}
    \varepsilon_{xx} & \varepsilon_{xy} & \varepsilon_{xz}\\
    \varepsilon_{xy} & \varepsilon_{yy} & \varepsilon_{yz}\\
    \varepsilon_{xz} & \varepsilon_{yz} & \varepsilon_{zz}
  \end{pmatrix}.
\]
Now $\vD$ need not point along $\vE$. The isotropic case is the diagonal one
with all three entries equal. The simulator treats the tensor case as the
general one and the scalar case as a special diagonal---a pattern we will see
repeatedly: handle the general object, specialize for free. The material
coefficient becomes a per-point weight $Q$ inside the bilinear form
$a(u,v)=(Q\,\diffop u,\diffop v)_\dom$ of \cref{ch:galerkin}, and whether $Q$
is a scalar or a $3\times3$ tensor changes nothing in the structure of the
assembly---only the arithmetic at each quadrature point.

\section{Interfaces and walls: boundary conditions}
\label{sec:bc}

A simulation lives in a bounded box, and the box is made of materials that
meet at interfaces and terminate at walls. What the fields do at those
surfaces is dictated by Maxwell's equations themselves, applied across the
discontinuity. Two conditions matter, and they are the continuous ancestors of
the \emph{essential} and \emph{natural} boundary conditions that the
finite-element method will later impose (\cref{ch:weak}, \cref{ch:bc}).

\begin{figure}[t]
\centering
\begin{tikzpicture}[scale=1.0]
  % two media separated by an interface
  \fill[simBgBlue]  (-3,0) rectangle (3,1.8);
  \fill[simBgGold]  (-3,-1.8) rectangle (3,0);
  \draw[very thick,simInk] (-3,0)--(3,0);
  \node[font=\footnotesize] at (-2.5,1.4) {medium 1};
  \node[font=\scriptsize,text=simGray] at (-2.5,1.0) {$\varepsilon_1,\mu_1$};
  \node[font=\footnotesize] at (-2.5,-1.4) {medium 2};
  \node[font=\scriptsize,text=simGray] at (-2.5,-1.0) {$\varepsilon_2,\mu_2$};
  % interface normal
  \draw[-{Stealth[length=2.2mm]},simRust,thick] (0,0)--(0,0.9)
     node[right,font=\scriptsize]{$\vn$};
  % tangential E continuous: same arrow both sides
  \draw[-{Stealth[length=2mm]},simBlue,thick] (1.2,0.0)--(2.2,0.0)
     node[above,font=\scriptsize]{$\vE_{\parallel}$};
  \draw[densely dashed,simBlue] (1.7,0.06)--(1.7,-0.5);
  \draw[-{Stealth[length=2mm]},simBlue,thick] (1.2,-0.0)--(2.2,-0.0);
  \node[font=\scriptsize,text=simBlue,align=center] at (1.7,-0.85)
     {$\vE_\parallel$ continuous};
  % normal B continuous: vertical arrow crossing
  \draw[-{Stealth[length=2.4mm]},simTeal,thick] (-1.7,-0.85)--(-1.7,0.85);
  \node[font=\scriptsize,text=simTeal] at (-1.15,0.55) {$\vB_\perp$};
  \node[font=\scriptsize,text=simTeal,align=center] at (-1.6,-1.15)
     {$\vB_\perp$ continuous};
\end{tikzpicture}
\caption[Interface boundary conditions]{At an interface between two media the
\emph{tangential} electric field $\vE_\parallel$ and the \emph{normal} magnetic
flux $\vB_\perp$ pass through unchanged. These two continuities are what the
finite-element spaces $\Hcurl$ (tangentially continuous, for $\vE$) and $\Hdiv$
(normally continuous, for $\vB$) are built to honour exactly. The complementary
components ($\vD_\perp$, $\vH_\parallel$) may jump by surface charge or current.}
\label{fig:interface}
\end{figure}

\subsection{Tangential \texorpdfstring{$\vE$}{E} and normal \texorpdfstring{$\vB$}{B} are continuous}

Apply Faraday's law to a thin loop straddling the interface, then shrink the
loop's height to zero. The flux through the vanishing loop goes to zero, and
what survives is the statement that the \emph{tangential component of $\vE$ is
continuous across the interface}:
\begin{equation}
  \vn\times(\vE_1-\vE_2) = \bm 0. \label{eq:bc-Etan}
\end{equation}
Symmetrically, apply Gauss's law for $\vB$ to a thin pillbox and shrink its
height: the \emph{normal component of $\vB$ is continuous},
\begin{equation}
  \vn\cdot(\vB_1-\vB_2) = 0. \label{eq:bc-Bnorm}
\end{equation}
The complementary components are \emph{not} continuous in general: the normal
$\vD$ can jump by a surface charge density, and the tangential $\vH$ can jump by
a surface current. But \eqref{eq:bc-Etan} and \eqref{eq:bc-Bnorm}---tangential
$\vE$, normal $\vB$---always hold. They are the reason the discrete spaces of
\cref{ch:functionspaces} are designed as they are: $\Hcurl$ enforces exactly
tangential continuity (the right home for $\vE$), and $\Hdiv$ enforces exactly
normal continuity (the right home for $\vB$). The physics chooses the function
space.

\subsection{PEC and PMC walls}

Two idealized walls bound nearly every simulation.

A \emph{perfect electric conductor} (PEC) is the $\sigma\to\infty$ limit of a
metal. Inside it the field is zero, so by tangential continuity
\eqref{eq:bc-Etan} the tangential electric field must vanish \emph{at the
surface}:
\begin{equation}
  \vn\times\vE = \bm 0 \quad\text{on a PEC wall.} \label{eq:pec}
\end{equation}
A PEC wall is a mirror that shorts out any tangential $\vE$; it is the
boundary condition at the metal plates of a capacitor and the walls of a
cavity. It will become an \emph{essential} (Dirichlet) condition: imposed
directly on the solution by removing the constrained degrees of freedom.

A \emph{perfect magnetic conductor} (PMC) is the dual idealization,
$\vn\times\vH=\bm0$. No real material is a PMC, but the condition arises
naturally as a \emph{symmetry plane}---a surface across which the problem is
mirror-symmetric---and it is the condition the finite-element method imposes by
doing \emph{nothing} (a \emph{natural} condition). The contrast is worth
holding onto: PEC must be enforced; PMC happens for free. That asymmetry,
traced to the integration by parts of \cref{ch:weak}, is why one of them costs
work to impose and the other costs nothing.

\begin{notationbox}
\textbf{Essential} (Dirichlet) conditions constrain the solution value and must
be \emph{built into} the function space---the simulator eliminates the
constrained degrees of freedom (\code{eliminate\_essential\_bc} in the
machinery of \cref{ch:bc}). \textbf{Natural} (Neumann-like) conditions
constrain a derivative and fall out of the weak form automatically; the
simulator imposes them by adding nothing. PEC $\to$ essential; PMC and
symmetry $\to$ natural. This single distinction reappears in every driver.
\end{notationbox}

\section{Energy and Poynting's theorem}
\label{sec:energy}

A simulation that finds the fields is only halfway done; the engineer wants a
\emph{number}---a capacitance, an inductance, a quality factor---and almost
every such number is an \emph{energy}. The electromagnetic field stores energy
in space at a density
\begin{equation}
  u = \tfrac12\,\vE\cdot\vD + \tfrac12\,\vH\cdot\vB
    = \tfrac12\,\varepsilon\,\abs{\vE}^2 + \tfrac{1}{2\mu}\,\abs{\vB}^2,
  \label{eq:energy-density}
\end{equation}
the first term electric, the second magnetic. Integrate $u$ over the device and
you have the total stored energy $U$---and from $U$ come the reductions:
capacitance is $C = 2U/V^2$ at unit voltage, inductance is $L = 2U/I^2$ at unit
current. These are the very quantities the \emph{energy reductions} of
\cref{ch:reductions} compute. Equation~\eqref{eq:energy-density} is their
integrand.

\subsection{Poynting's theorem: where the energy goes}

Energy density alone is a snapshot; we also want to know how energy
\emph{moves}. That is the content of Poynting's theorem, and it follows from
Maxwell's equations by one inspired manipulation. Define the \emph{Poynting
vector}
\begin{equation}
  \vS = \vE\times\vH, \label{eq:poynting-vec}
\end{equation}
the directional flow of electromagnetic power per unit area (\cref{fig:poynting}).
Now form the combination $\vH\cdot(\Curl\vE) - \vE\cdot(\Curl\vH)$. A standard
vector identity rewrites it as a divergence,
\[
  \vH\cdot(\Curl\vE) - \vE\cdot(\Curl\vH) = \Div(\vE\times\vH) = \Div\vS.
\]
Substitute Faraday's law~\eqref{eq:faraday} for $\Curl\vE$ and the
Amp\`ere--Maxwell law~\eqref{eq:ampere} for $\Curl\vH$:
\[
  \Div\vS
  = \vH\cdot(-\partial_t\vB) - \vE\cdot(\vJ + \partial_t\vD)
  = -\,\partial_t\!\left(\tfrac12\vH\cdot\vB + \tfrac12\vE\cdot\vD\right)
    - \vE\cdot\vJ.
\]
Recognizing the bracket as the energy density $u$ of
\eqref{eq:energy-density}, we arrive at \emph{Poynting's theorem}:
\begin{equation}
  \partial_t u + \Div\vS = -\,\vE\cdot\vJ. \label{eq:poynting}
\end{equation}
Read it as a balance sheet for energy at every point. The first term is the
rate the stored energy changes; the second is the net power flowing \emph{out}
(the divergence of the energy flux $\vS$); the right side is the power
delivered to the currents. When $\vJ=\sigma\vE$ (Ohmic material), the source
term is $-\sigma\abs{\vE}^2\le 0$: energy is \emph{dissipated}, drained
irreversibly into heat. That loss term is what gives a resonator its finite
$Q$.

\begin{figure}[t]
\centering
\begin{tikzpicture}[scale=1.0]
  % E field (blue, right), H field (green, up/out), S = E x H (rust, into-page-ish)
  \draw[-{Stealth[length=2.6mm]},simBlue,very thick] (0,0)--(2.0,0)
     node[right,font=\small]{$\vE$};
  \draw[-{Stealth[length=2.6mm]},simGreen,very thick] (0,0)--(0,2.0)
     node[above,font=\small]{$\vH$};
  \draw[-{Stealth[length=2.6mm]},simRust,very thick] (0,0)--(1.45,1.45)
     node[above right,font=\small]{$\vS=\vE\times\vH$};
  % little flux disk
  \fill[simRust!18] (0.0,0.0) -- (2.0,0.0) arc (0:90:2.0) -- cycle;
  \node[font=\scriptsize,text=simGray,align=center] at (3.6,1.0)
     {power flows\\ along $\vS$};
  % right angle ticks
  \draw[simInk] (0.18,0)--(0.18,0.18)--(0,0.18);
\end{tikzpicture}
\caption[The Poynting vector]{The Poynting vector $\vS=\vE\times\vH$ points
perpendicular to both fields and gives the direction and density of
electromagnetic power flow. In a propagating wave, $\vE$, $\vH$, and $\vS$ form
a right-handed triple and $\vS$ points the way the wave travels. Poynting's
theorem \eqref{eq:poynting} is the local accounting that $\partial_t u+\Div\vS$
equals the power delivered to charges.}
\label{fig:poynting}
\end{figure}

\begin{keyidea}
Poynting's theorem is the conservation of energy for the field, derived
\emph{from} Maxwell's equations rather than assumed alongside them---just as
continuity \eqref{eq:continuity} was conservation of charge. The pattern is
worth noticing: the deep facts of the theory (energy conservation, charge
conservation) are not extra axioms but algebraic consequences of the four laws,
extracted by a vector identity. The energy density \eqref{eq:energy-density} is
the integrand the output reductions sum; the loss term $\vE\cdot\vJ$ is what
makes a quality factor finite.
\end{keyidea}

\subsection{Stored energy, the quality factor, and participation}
\label{sec:Q}

The energy bookkeeping is not an afterthought; it is what most analyses
\emph{report}. Three quantities, all built from the energy density
\eqref{eq:energy-density}, recur throughout the book.

First, in a resonant mode the electric and magnetic energies are equal on
average and exchange completely twice per cycle---energy sloshes from the
electric field to the magnetic field and back, just as a pendulum trades
kinetic for potential energy. The total $U=U_E+U_B$ is what the eigenmode
analysis integrates.

Second, when a small loss is present, the mode does not ring forever; its
amplitude decays, and the \emph{quality factor}
\[
  Q = \omega\,\frac{U_{\text{stored}}}{P_{\text{loss}}}
    = \frac{\omega}{\kappa}
\]
measures how many radians of oscillation it takes to drain a factor $1/e$ of the
energy. A high-$Q$ cavity ($Q\sim10^6$ or more for a superconducting resonator)
rings for millions of cycles; a lossy one barely rings at all. The numerator is
the stored field energy of \eqref{eq:energy-density}; the denominator is the
power-loss integral of $\sigma\abs{\vE}^2$ from Poynting's
theorem~\eqref{eq:poynting}. This is the $Q$ the eigenmode reduction of
\cref{ch:reductions} computes, and the reason that chapter cares about the loss
term at all.

Third, engineers want to know \emph{where} a mode keeps its energy---in which
dielectric, in which gap. The \emph{participation ratio} of a region is the
fraction of the total energy stored there,
\[
  p_{\text{region}} = \frac{\int_{\text{region}} u\dV}{\int_{\text{all}} u\dV},
\]
a number between $0$ and $1$ that sums to one over a partition of the device. It
is the headline output of a great deal of quantum-device design: the
participation of a lossy interface sets the limit on a qubit's coherence. Like
$Q$ and like capacitance, it is just the energy density
\eqref{eq:energy-density} integrated and divided---which is why a single
\emph{energy reduction} (\cref{ch:reductions}) computes all three.

\section{Three regimes: static, time-harmonic, time-domain}
\label{sec:regimes}

The full Maxwell system is a coupled set of time-dependent partial differential
equations. Almost no engineering question needs it in that generality. Instead
each analysis lives in one of three \emph{regimes}, distinguished entirely by
\emph{how time enters}, and the choice of regime is what turns the one physics
into the several analyses of \cref{ch:targets}. \Cref{fig:regimes} is the map;
the rest of this section walks it.

\begin{figure}[t]
\centering
\begin{tikzpicture}[node distance=8mm]
  \node[stage, minimum width=30mm, minimum height=14mm] (full)
    {\textbf{Maxwell}\\[1pt]\footnotesize $\partial_t$ present};
  \node[stage, minimum width=30mm, below left=14mm and 8mm of full] (static)
    {\textbf{Static}\\[1pt]\footnotesize $\partial_t\to 0$};
  \node[stage, minimum width=30mm, below=22mm of full] (harm)
    {\textbf{Time-harmonic}\\[1pt]\footnotesize $\partial_t\to i\omega$};
  \node[stage, minimum width=30mm, below right=14mm and 8mm of full] (td)
    {\textbf{Time-domain}\\[1pt]\footnotesize $\partial_t$ kept};
  \draw[flow] (full) -- (static);
  \draw[flow] (full) -- (harm);
  \draw[flow] (full) -- (td);
  % products
  \node[product, below=7mm of static, minimum width=30mm, font=\scriptsize]
    {electro- / magnetostatic\\ $\to$ $C$, $L$};
  \node[product, below=7mm of harm, minimum width=34mm, font=\scriptsize]
    {eigenmode, driven\\ $\to$ $f,Q$, S-params};
  \node[product, below=7mm of td, minimum width=30mm, font=\scriptsize]
    {transient\\ $\to$ field history};
  % the key substitution callout
  \node[font=\scriptsize\itshape, text=simRust, right=10mm of harm, align=center]
    (sub) {the phasor\\ substitution};
  \draw[backflow] (sub) -- (harm);
\end{tikzpicture}
\caption[The three regimes]{The three regimes of Maxwell's equations, set by
how the time derivative is treated. Dropping $\partial_t$ gives the
\emph{static} problems (electrostatic, magnetostatic). Replacing
$\partial_t\to i\omega$ gives the \emph{time-harmonic} problems (eigenmode,
driven) and produces the complex curl--curl operators of
\cref{ch:reductions-pde}. Keeping $\partial_t$ gives the \emph{time-domain}
transient problem. The single substitution $\partial_t\to i\omega$ is the
bridge from the time-domain to the frequency-domain analyses.}
\label{fig:regimes}
\end{figure}

\subsection{Static: time stands still}

When nothing changes in time, every $\partial_t$ vanishes and the two coupling
laws \emph{decouple}. Faraday's law becomes $\Curl\vE=\bm0$, so $\vE$ is a
gradient (\cref{sec:potentials}); Amp\`ere's law loses its displacement current
and becomes $\Curl\vH=\vJ$. The electric and magnetic problems no longer talk
to each other, and we are left with two independent, self-contained problems:

\begin{itemize}
  \item \textbf{Electrostatics.} With $\vE=-\Grad V$ and $\vD=\varepsilon\vE$,
  Gauss's law $\Div\vD=\rho$ becomes a single scalar second-order equation for
  the potential,
  \begin{equation}
    -\Div(\varepsilon\,\Grad V) = \rho, \label{eq:poisson}
  \end{equation}
  the \emph{Poisson equation} (Laplace's equation $-\Div(\varepsilon\Grad
  V)=0$ in charge-free regions). This is the electrostatic analysis, and it is
  the simplest governing equation in the book.

  \item \textbf{Magnetostatics.} With $\vB=\Curl\vA$ and $\vH=\nu\vB$,
  Amp\`ere's law $\Curl\vH=\vJ$ becomes the \emph{curl--curl} equation for the
  vector potential,
  \begin{equation}
    \Curl(\nu\,\Curl\vA) = \vJ. \label{eq:curlcurl-static}
  \end{equation}
\end{itemize}

These are the equations behind the electrostatic and magnetostatic analyses of
\cref{ch:targets}, and \cref{ch:reductions-pde} derives them in full. Notice
already how alike they are in shape---a weighted second-order operator equal to
a source---differing only in the function space ($V$ scalar, $\vA$ vector) and
the differential operator ($\Grad$ versus $\Curl$). That kinship is the seed of
``the same machine, one part swapped.''

\subsection{Time-harmonic: the phasor substitution}

Most microwave engineering asks not ``what happens in time?'' but ``what is the
steady response at frequency $\omega$?'' If every field oscillates at a single
angular frequency $\omega$, we can write each as the real part of a complex
\emph{phasor} times $e^{i\omega t}$:
\[
  \vE(\vx,t) = \Re\!\big\{\hat{\vE}(\vx)\,e^{i\omega t}\big\},
\]
and similarly for $\vB,\vD,\vH,\vJ$. The complex amplitude $\hat{\vE}(\vx)$
carries both magnitude and phase at every point; time has been factored out.
The payoff is a single, beautiful simplification. Differentiating
$e^{i\omega t}$ in time just multiplies by $i\omega$:
\begin{equation}
  \partial_t \;\longrightarrow\; i\omega. \label{eq:phasor-sub}
\end{equation}
Every time derivative in Maxwell's equations becomes an algebraic factor
$i\omega$, and the partial differential equations \emph{in space and time}
collapse to partial differential equations \emph{in space alone}, with complex
coefficients. The phasor curl laws read
\begin{align}
  \Curl\hat{\vE} &= -\,i\omega\,\hat{\vB}, \label{eq:faraday-phasor}\\
  \Curl\hat{\vH} &= \hat{\vJ} + i\omega\,\hat{\vD}. \label{eq:ampere-phasor}
\end{align}
Worked example~\ref{wex:phasor} carries this substitution all the way to the curl--curl
operator that the driven and eigenmode analyses solve. The substitution
\eqref{eq:phasor-sub} is the single most important bridge in Part~II: it is
what turns the time-domain Maxwell system into the frequency-domain operators
$\mat K + i\omega\mat C - \omega^2\mat M$ that \cref{ch:reductions-pde} derives
and that \cref{part:machinery} learns to solve.

It is worth seeing in advance where those three operators come from, because the
pattern recurs in every harmonic analysis. After the weak form and
discretization of \cref{ch:weak,ch:galerkin}, each \emph{term} of the harmonic
curl--curl equation contributes one matrix:
\begin{itemize}
  \item the curl--curl term $\Curl(\nu\,\Curl\hat{\vE})$ becomes the
  \emph{stiffness} $\mat K$ (no $\omega$ in it---it is the static operator);
  \item a conduction/loss term $i\omega\sigma\hat{\vE}$ (present when the medium
  conducts, or a port radiates) becomes $i\omega\mat C$, the \emph{damping}
  operator---first power of $\omega$, and the only one carrying the imaginary
  unit;
  \item the displacement term $-\omega^2\varepsilon\hat{\vE}$ becomes
  $-\omega^2\mat M$, with the \emph{mass} $\mat M$---second power of $\omega$,
  the term that makes the problem a resonance.
\end{itemize}
The full operator $\mat A(\omega)=\mat K + i\omega\mat C - \omega^2\mat M$ is a
\emph{quadratic} polynomial in $\omega$ whose three coefficient matrices are
assembled once and reused at every frequency. Drop the drive and you have the
eigenmode problem $\mat K\hat{\vE}=\omega^2\mat M\hat{\vE}$ (a generalized
eigenproblem when $\mat C=0$, a quadratic one when it does not); keep the drive
and you sweep $\omega$, re-forming $\mat A(\omega)$ and solving
$\mat A(\omega)\hat{\vE}=\vb(\omega)$ at each step. The single phasor
substitution has produced the entire frequency-domain operator algebra the
later chapters live in.

\subsection{Time-domain: keep the clock running}

Sometimes the excitation is a broadband pulse, or the materials are nonlinear,
or one simply wants the movie. Then no single frequency suffices and the time
derivatives must be \emph{kept}. After spatial discretization the result is a
second-order system of ordinary differential equations in time,
\[
  \mat M\,\ddot{\vs}(t) + \mat C\,\dot{\vs}(t) + \mat K\,\vs(t) = \vJ(t),
\]
the same $\mat M,\mat C,\mat K$ operators as the harmonic case but now marched
forward step by step rather than solved frequency by frequency. This is the
transient analysis. Its defining feature---each step begins from the state the
previous step produced---makes it the book's first genuinely \emph{sequential}
computation, and the motivating example for the fold of \cref{ch:fold}.

\begin{keyidea}
One physics, three regimes, set by the fate of $\partial_t$. \textbf{Drop} it
($\partial_t\to0$): static, real operators, two decoupled problems (electro-
and magnetostatic). \textbf{Algebraize} it ($\partial_t\to i\omega$):
time-harmonic, complex operators, the curl--curl/Helmholtz family (eigenmode,
driven). \textbf{Keep} it: time-domain, a marched ODE system (transient). The
substitution $\partial_t\to i\omega$ is the hinge between the second and third,
and the source of every complex operator in the book.
\end{keyidea}

\begin{table}[t]
\centering
\caption[The three regimes, compared]{The three regimes side by side. The
column that changes everything is the fate of the time derivative; it dictates
whether the operator is real or complex, whether the problem couples $\vE$ and
$\vB$, and which analyses of \cref{ch:targets} it produces. The discrete
operators $\mat K,\mat C,\mat M$ are the same matrices in every column---only
their combination differs.}
\label{tab:regimes}
\small
\setlength{\tabcolsep}{5pt}
\renewcommand{\arraystretch}{1.3}
\begin{tabular}{@{}llll@{}}
\toprule
 & Static & Time-harmonic & Time-domain \\
\midrule
$\partial_t$ becomes & $0$ & $i\omega$ (algebraic) & kept (a derivative) \\
Unknown & real field & complex phasor & real field history \\
Coupling & decoupled & coupled & coupled \\
Discrete operator & $\mat K$ & $\mat K+i\omega\mat C-\omega^2\mat M$
  & $\mat M\ddot{\vs}+\mat C\dot{\vs}+\mat K\vs$ \\
Solve shape & one linear solve & solve per $\omega$ / eigenproblem & marched fold \\
Analyses & electro-/magnetostatic & eigenmode, driven & transient \\
\bottomrule
\end{tabular}
\end{table}

\section{Potentials: solving the constraints automatically}
\label{sec:potentials}

The static analyses were posed not for the fields $\vE,\vB$ directly but for
\emph{potentials}---a scalar $V$ and a vector $\vA$. The reason is elegant: a
potential \emph{automatically satisfies} one of the Gauss laws, removing it
from the list of equations to solve. Potentials trade a constrained field for an
unconstrained one.

\subsection{The scalar potential \texorpdfstring{$V$}{V}}

In the static (or curl-free) case Faraday's law gives $\Curl\vE=\bm0$. A vector
identity from vector calculus states that the curl of \emph{any} gradient
vanishes, $\Curl(\Grad f)=\bm0$. So if we write
\begin{equation}
  \vE = -\Grad V, \label{eq:E-from-V}
\end{equation}
then $\Curl\vE = -\Curl(\Grad V)=\bm 0$ is satisfied \emph{identically}, for
\emph{any} scalar function $V$ (\cref{fig:potentials}, left). The curl law is no
longer an equation to enforce; it is built into the representation. We are left
with only Gauss's law, which \eqref{eq:E-from-V} turns into the Poisson
equation~\eqref{eq:poisson}. The minus sign is convention: it makes $\vE$ point
``downhill,'' from high potential to low, matching the picture of a positive
charge rolling down a voltage slope.

\subsection{The vector potential \texorpdfstring{$\vA$}{A}}

The magnetic side runs in parallel. Gauss's law for $\vB$ says $\Div\vB=0$.
The companion identity is that the divergence of \emph{any} curl vanishes,
$\Div(\Curl\vF)=0$ (the same identity that gave us continuity in
\eqref{eq:continuity}). So if we write
\begin{equation}
  \vB = \Curl\vA, \label{eq:B-from-A}
\end{equation}
then $\Div\vB=\Div(\Curl\vA)=0$ holds \emph{identically} for any vector field
$\vA$ (\cref{fig:potentials}, right). Gauss's law for $\vB$ disappears, and
Amp\`ere's law becomes the curl--curl equation
\eqref{eq:curlcurl-static}. The two potentials are the two halves of the same
trick: each rides on a vanishing-derivative identity ($\Curl\Grad=0$ and
$\Div\Curl=0$) to dissolve one Gauss law and leave one equation behind. Those
two identities are exactly the two composites of the de Rham complex
$\Hone\xrightarrow{\Grad}\Hcurl\xrightarrow{\Curl}\Hdiv\xrightarrow{\Div}\Ltwo$
that \cref{ch:derham} builds the discrete function spaces to honour.

\begin{figure}[t]
\centering
\begin{tikzpicture}[scale=1.0]
% left: scalar potential, E = -grad V over a hill
\begin{scope}[shift={(0,0)}]
  % contour "hill"
  \foreach \r/\op in {1.05/12, 0.75/22, 0.45/36}{
    \draw[simGray!60] (0,0) circle (\r);}
  \node[font=\scriptsize] at (0,0) {high $V$};
  % E arrows pointing downhill (outward, but E=-grad V so toward low V outside)
  \foreach \a in {30,120,210,300}{
    \draw[-{Stealth[length=2mm]},simBlue,thick]
      (\a:0.5)--(\a:1.25);}
  \node[font=\scriptsize\itshape,simBlue] at (1.35,1.0) {$\vE=-\Grad V$};
  \node[font=\footnotesize] at (0,-1.7) {scalar potential};
\end{scope}
% right: vector potential, B = curl A
\begin{scope}[shift={(5.0,0)}]
  % A as a circulation
  \draw[flow,simViolet] (0,0) circle (0.8);
  \node[font=\scriptsize\itshape,simViolet] at (-1.05,0.6) {$\vA$};
  % B coming out of the loop
  \draw[-{Stealth[length=2.4mm]},simTeal,very thick] (0,0)--(0,0.0);
  \fill[simTeal] (0,0) circle (1.6pt);
  \node[font=\scriptsize\itshape,simTeal] at (0.45,0.25) {$\vB$};
  \node[font=\footnotesize] at (0,-1.7) {vector potential};
  \node[font=\scriptsize,text=simGray] at (0,-1.4) {$\vB=\Curl\vA$};
\end{scope}
\end{tikzpicture}
\caption[The potential representations]{The two potential representations.
\emph{Left}: writing $\vE=-\Grad V$ makes $\vE$ point downhill on the potential
landscape and satisfies $\Curl\vE=\bm0$ identically. \emph{Right}: writing
$\vB=\Curl\vA$ makes $\vB$ the circulation of $\vA$ and satisfies $\Div\vB=0$
identically. Each potential dissolves one Gauss law for free, riding on a
vanishing-derivative identity of vector calculus.}
\label{fig:potentials}
\end{figure}

\subsection{Gauge freedom, briefly}

There is a price, gentle but real. The potentials are not unique. Adding any
constant to $V$ leaves $\vE=-\Grad V$ unchanged. Adding any gradient $\Grad
\chi$ to $\vA$ leaves $\vB=\Curl\vA$ unchanged, because $\Curl(\Grad\chi)=\bm0$.
This freedom---the \emph{gauge freedom}---means the magnetostatic curl--curl
operator \eqref{eq:curlcurl-static} has a nontrivial null space: every gradient
field is invisible to it. One fixes the freedom by an extra condition (a
\emph{gauge}, such as $\Div\vA=0$), or, as the simulator does, by projecting
the solution onto the gradient-free part with a divergence-free projector
(\cref{ch:bc}). For now simply note that the null space is not a bug to be
feared but the direct fingerprint of gauge freedom---the same vanishing-curl
identity that made the potential useful, seen from the other side.

\section{Worked examples}

The examples below are deliberately small enough to check with a pencil. The
first exercises Gauss's law and the energy integral end to end---it is the
electrostatic analysis in miniature. The second carries the phasor substitution
to the operator the time-harmonic analyses actually solve. The third puts the
interface conditions of \cref{sec:bc} to work on a concrete refraction of field
lines.

\begin{workedexample}{Parallel-plate capacitor: field, charge, energy}{capacitor}
Two parallel conducting plates of area $A$ sit a distance $d$ apart, the gap
filled with a dielectric of permittivity $\varepsilon=\varepsilon_r\varepsilon_0$.
The top plate is held at potential $V_0$, the bottom at $0$ (\cref{fig:capacitor}).
Find the field, the charge, the capacitance, and the stored energy---ignoring
fringing, so the field is uniform and vertical between the plates.

\emph{Field from the potential.} With the plates close compared to their size,
the potential varies linearly across the gap: $V(z)=V_0\,z/d$ for $0\le z\le
d$. Then
\[
  \vE = -\Grad V = -\frac{dV}{dz}\,\hat{\vz} = -\frac{V_0}{d}\,\hat{\vz},
\]
a uniform field of magnitude $E=V_0/d$ pointing from the high plate to the low.
With $V_0=\SI{1}{\volt}$ and $d=\SI{1}{\micro\meter}$ this is
$E=\SI{e6}{\volt\per\meter}$---fields inside small devices are large.

\emph{Charge from Gauss's law.} The displacement is $\vD=\varepsilon\vE$, of
magnitude $D=\varepsilon V_0/d$. A pillbox Gauss surface straddling the top
plate sees flux $D$ leave through its face of area $A$, so the enclosed surface
charge is
\[
  Q = D\,A = \frac{\varepsilon A}{d}\,V_0.
\]

\emph{Capacitance.} By definition $C=Q/V_0$, so
\begin{equation}
  C = \frac{\varepsilon A}{d}. \label{eq:cap}
\end{equation}
The classic result: capacitance grows with plate area and dielectric constant,
and shrinks with separation.

\emph{Energy, two ways.} The stored energy from the energy
density~\eqref{eq:energy-density}, integrated over the gap volume $V_{\text{gap}}=Ad$,
is
\[
  U = \int \tfrac12\varepsilon\abs{\vE}^2\dV
    = \tfrac12\varepsilon\Big(\frac{V_0}{d}\Big)^{\!2}\,(Ad)
    = \tfrac12\,\frac{\varepsilon A}{d}\,V_0^2
    = \tfrac12\,C\,V_0^2,
\]
recovering the familiar $U=\tfrac12 C V_0^2$. Read backwards, this says
$C=2U/V_0^2$: \emph{the capacitance is twice the field energy at unit
voltage}---exactly the energy reduction the electrostatic analysis performs in
\cref{ch:reductions}. A concrete number: $\varepsilon_r=3.9$ (silicon dioxide),
$A=\SI{100}{\micro\meter\squared}=\SI{e-10}{\meter\squared}$,
$d=\SI{100}{\nano\meter}$ gives
$C=3.9\times\SI{8.854e-12}{}\times\SI{e-10}{}/\SI{e-7}{}\approx\SI{3.5e-14}{\farad}=\SI{35}{\femto\farad}$,
a typical on-chip value.
\end{workedexample}

\begin{figure}[t]
\centering
\begin{physicsbox}
\physicsframe{
  \fill[simBgBlue] (-1.6,-0.9) rectangle (1.6,0.9);
  \draw[very thick,simInk] (-1.6,0.9)--(1.6,0.9);
  \draw[very thick,simInk] (-1.6,-0.9)--(1.6,-0.9);
  \node[font=\scriptsize] at (-2.0,0.9) {$V_0$};
  \node[font=\scriptsize] at (-2.0,-0.9) {$0$};
  \node[font=\scriptsize,text=simGray] at (0,1.25) {area $A$};
  \draw[{Stealth[length=1.5mm]}-{Stealth[length=1.5mm]},simGray]
    (1.95,0.9)--(1.95,-0.9);
  \node[font=\scriptsize,text=simGray] at (2.25,0) {$d$};
  \foreach \x in {-1.2,-0.6,0,0.6,1.2}{
    \draw[-{Stealth[length=1.8mm]},simBlue,thick] (\x,0.82)--(\x,-0.82);}
  \node[font=\scriptsize\itshape,simBlue] at (1.0,0.45) {$\vE$};
  \node[font=\scriptsize,text=simGreen] at (0,0) {$\varepsilon$};
}
\end{physicsbox}
\caption[Parallel-plate capacitor]{The parallel-plate capacitor of
Worked example~\ref{wex:capacitor}. A uniform vertical field $\vE=-(V_0/d)\hat{\vz}$ fills the
dielectric-filled gap; Gauss's law gives the surface charge, and the field
energy gives the capacitance $C=\varepsilon A/d = 2U/V_0^2$.}
\label{fig:capacitor}
\end{figure}

\begin{workedexample}{Phasor reduction to the curl--curl operator}{phasor}
Take a source-free region ($\vJ=\bm0$, $\rho=0$) of a linear medium
($\vD=\varepsilon\vE$, $\vB=\mu\vH$). Show that the time-harmonic Maxwell
equations reduce to a single second-order operator equation for $\hat{\vE}$,
the \emph{curl--curl} (vector Helmholtz) equation that the driven and eigenmode
analyses solve.

\emph{Start from the phasor curl laws.} With $\vJ=\bm0$, equations
\eqref{eq:faraday-phasor}--\eqref{eq:ampere-phasor} read
\[
  \Curl\hat{\vE} = -\,i\omega\,\hat{\vB} = -\,i\omega\mu\,\hat{\vH},
  \qquad
  \Curl\hat{\vH} = i\omega\,\hat{\vD} = i\omega\varepsilon\,\hat{\vE}.
\]

\emph{Eliminate $\hat{\vH}$.} Solve the first for $\hat{\vH}=
-\,\dfrac{1}{i\omega\mu}\Curl\hat{\vE}$ and substitute into the second:
\[
  \Curl\!\left(-\frac{1}{i\omega\mu}\Curl\hat{\vE}\right)
    = i\omega\varepsilon\,\hat{\vE}.
\]
Multiply through by $i\omega$ and write $\nu=1/\mu$. The factor $i\omega$ cancels
on the left to leave $-\Curl(\nu\Curl\hat{\vE})$ on one side, and $i\omega\cdot
i\omega = -\omega^2$ on the other:
\begin{equation}
  \Curl(\nu\,\Curl\hat{\vE}) - \omega^2\varepsilon\,\hat{\vE} = \bm0.
  \label{eq:curlcurl-harmonic}
\end{equation}

\emph{Read the result.} Equation~\eqref{eq:curlcurl-harmonic} is the
\emph{curl--curl} operator equation. Its two terms are exactly the discrete
operators of \cref{ch:targets}: the curl--curl term becomes the stiffness
$\mat K$ (weighted by reluctivity $\nu$), and the $\omega^2$ term becomes the
mass $\mat M$ (weighted by permittivity $\varepsilon$). With no drive it is the
generalized eigenvalue problem $\mat K\hat{\vE}=\omega^2\mat M\hat{\vE}$---the
\emph{eigenmode} analysis, whose eigenvalues are the squared resonant
frequencies. Restore a current source $\hat{\vJ}$ and add a loss term
$i\omega\sigma$ (from $\vJ=\sigma\vE$) and the operator becomes
$\mat K+i\omega\mat C-\omega^2\mat M$---the \emph{driven} analysis. Both
analyses are this one phasor reduction, with and without a right-hand side.
This is the single calculation that turns the four time-domain laws into the
frequency-domain machinery the rest of the book solves.
\end{workedexample}

\begin{workedexample}{Refraction of electric field lines at an interface}{refract}
A static electric field crosses the planar interface between two dielectrics,
$\varepsilon_1$ above and $\varepsilon_2$ below, with no free surface charge.
On the upper side the field makes an angle $\theta_1$ with the normal. Find the
angle $\theta_2$ on the lower side, and show that field lines \emph{bend} at the
interface like light through glass.

\emph{Set up the two conditions.} Resolve each field into a normal component
(along $\vn$) and a tangential component (in the surface). The interface
conditions of \cref{sec:bc} give two equations. Tangential $\vE$ is continuous
\eqref{eq:bc-Etan}:
\[
  E_1\sin\theta_1 = E_2\sin\theta_2.
\]
With no surface charge, normal $\vD=\varepsilon\vE$ is continuous:
\[
  \varepsilon_1 E_1\cos\theta_1 = \varepsilon_2 E_2\cos\theta_2.
\]

\emph{Divide to eliminate the magnitudes.} Dividing the first equation by the
second cancels $E_1$ and $E_2$ and leaves a clean relation between the angles:
\begin{equation}
  \frac{\tan\theta_1}{\varepsilon_1} = \frac{\tan\theta_2}{\varepsilon_2},
  \qquad\text{i.e.}\qquad
  \tan\theta_2 = \frac{\varepsilon_2}{\varepsilon_1}\,\tan\theta_1.
  \label{eq:refract}
\end{equation}
This is the electrostatic ``law of refraction.''

\emph{Put numbers in.} Take air above ($\varepsilon_1=\varepsilon_0$) and a
dielectric below with $\varepsilon_2=4\varepsilon_0$, and let the field enter at
$\theta_1=\SI{30}{\degree}$. Then $\tan\theta_2 = 4\tan\SI{30}{\degree} =
4\times 0.577 = 2.31$, so $\theta_2 = \arctan(2.31)\approx\SI{66.6}{\degree}$.
The field line bends \emph{away} from the normal on entering the denser
dielectric---the opposite sense from a light ray, because here it is the normal
\emph{flux} $\vD$, not the field $\vE$, that is conserved. The lesson for the
simulator is concrete: the two interface conditions of \cref{sec:bc} are not
decorative; they are algebraic constraints the discrete field must satisfy at
every material boundary, and they are exactly what the $\Hcurl$ and $\Hdiv$
function spaces of \cref{ch:functionspaces} are constructed to enforce.
\end{workedexample}

\begin{pitfall}
The cancellation in Worked example~\ref{wex:phasor} relied on $\omega\neq0$. At $\omega=0$ the
phasor substitution \eqref{eq:phasor-sub} sends $\partial_t\to0$ and we are back
in the static regime, where the mass term vanishes and the curl--curl equation
degenerates to the magnetostatic $\Curl(\nu\Curl\vA)=\vJ$. The three regimes are
not separate physics; the static case is the $\omega\to0$ limit of the harmonic
one. Watch the limit: the $\omega^2\mat M$ term that defines a resonance simply
switches off.
\end{pitfall}

\section{A third look: the plane wave}
\label{sec:planewave}

It is worth seeing the solution that the source-free curl--curl equation
\eqref{eq:curlcurl-harmonic} hides: the \emph{plane wave}, the simplest
nontrivial field in all of electromagnetism and the one that makes ``Maxwell's
equations are light'' concrete. In a uniform medium try
$\hat{\vE}=\vE_0\,e^{-i\vk\cdot\vx}$, a field of fixed amplitude $\vE_0$ and a
single propagation direction $\vk$ (\cref{fig:planewave}). The curl of such a
field turns into $-i\vk\times$, and \eqref{eq:curlcurl-harmonic} becomes the
algebraic \emph{dispersion relation}
\[
  \abs{\vk}^2 = \omega^2\mu\varepsilon,
  \qquad\text{i.e.}\qquad
  \abs{\vk} = \frac{\omega}{c}, \quad c=\frac{1}{\sqrt{\mu\varepsilon}},
\]
the statement that the wave travels at speed $c=1/\sqrt{\mu\varepsilon}$. In
vacuum $c=1/\sqrt{\mu_0\varepsilon_0}\approx\SI{3e8}{\meter\per\second}$---the
speed of light, recovered from $\mu_0$ and $\varepsilon_0$ alone, with no light
ever mentioned. Faraday's law \eqref{eq:faraday-phasor} then forces
$\vk\cdot\vE_0=0$ (the field is transverse to the direction of travel) and ties
$\vH_0$ perpendicular to both, so $\vE$, $\vH$, and the Poynting vector
$\vS=\vE\times\vH$ form the right-handed triple of \cref{fig:poynting}, with
$\vS$ pointing along $\vk$.

\begin{figure}[t]
\centering
\begin{tikzpicture}[scale=1.0]
  % propagation axis
  \draw[->,simGray] (-0.3,0)--(6.4,0) node[right,font=\scriptsize]{$\vk$};
  % E field: blue sinusoid in the vertical plane
  \draw[simBlue,thick,domain=0:5.9,samples=100,smooth]
    plot (\x,{0.85*sin(180*\x)});
  \node[font=\scriptsize\itshape,simBlue] at (0.55,1.05) {$\vE$};
  % a few E vectors riding the crest
  \foreach \x in {0.5,1.5,2.5,3.5,4.5,5.5}{
    \draw[-{Stealth[length=1.3mm]},simBlue,semithick]
      (\x,0)--(\x,{0.85*sin(180*\x)});}
  % H field: smaller transverse green sinusoid in phase
  \draw[simGreen,thick,domain=0:5.9,samples=100,smooth]
    plot (\x,{0.40*sin(180*\x)});
  \node[font=\scriptsize\itshape,simGreen] at (1.55,-0.62) {$\vH$};
  \node[font=\scriptsize,text=simRust] at (5.0,0.95) {$\vS\parallel\vk$};
\end{tikzpicture}
\caption[A plane wave]{A plane electromagnetic wave. The electric field $\vE$
(blue) and magnetic field $\vH$ (green, drawn here as the smaller transverse
component) oscillate in phase, perpendicular to each other and to the
propagation direction $\vk$. The wave moves at speed $c=1/\sqrt{\mu\varepsilon}$
set by the dispersion relation $\abs{\vk}=\omega/c$, carrying power along
$\vS=\vE\times\vH$. This solution falls out of the same source-free curl--curl
operator \eqref{eq:curlcurl-harmonic} the eigenmode and driven analyses
assemble.}
\label{fig:planewave}
\end{figure}

\subsection{Lossy media and the skin depth}

The dispersion relation already carries the materials inside it, and turning on
conductivity makes the wave \emph{decay}. In a conductor the
Amp\`ere--Maxwell law gains the conduction current $\vJ=\sigma\vE$, and the
phasor curl law becomes $\Curl\hat{\vH}=(\sigma+i\omega\varepsilon)\hat{\vE}$.
Carrying this through the elimination of Worked example~\ref{wex:phasor} replaces
$\omega^2\mu\varepsilon$ by $\omega^2\mu\varepsilon - i\omega\mu\sigma$, so the
wavenumber $\abs{\vk}$ acquires an imaginary part. A complex $\vk=\beta - i/\delta$
means the field $e^{-i\vk\cdot\vx}$ carries a real exponential decay
$e^{-x/\delta}$: the wave dies away as it penetrates, over a characteristic
length---the \emph{skin depth}
\[
  \delta = \sqrt{\frac{2}{\omega\mu\sigma}}\qquad(\text{good conductor}).
\]
At microwave frequencies $\delta$ in copper is well under a micron, which is why
fields do not penetrate metal and why the PEC idealization
$\vn\times\vE=\bm0$ of \cref{sec:bc} is so accurate: a real conductor confines
the field to a vanishingly thin surface layer. The same imaginary part is the
microscopic origin of the loss term $-\sigma\abs{\vE}^2$ in Poynting's theorem
\eqref{eq:poynting}, and of the finite quality factor $Q$ a resonant cavity
exhibits.

Light is Maxwell's equations solving themselves in empty space; a damped wave
in a conductor is the same equations solving themselves in a lossy one. That
both fall out of the single curl--curl operator the simulator assembles---the
operator Worked example~\ref{wex:phasor} derived---is the unity the whole book is about.

\summarysection
Maxwell's four equations split into two \emph{Gauss} constraints
($\Div\vD=\rho$, $\Div\vB=0$) and two \emph{curl} dynamics (Faraday
$\Curl\vE=-\partial_t\vB$, Amp\`ere--Maxwell $\Curl\vH=\vJ+\partial_t\vD$); the
displacement current closes the loop that becomes light, and forces the
\emph{continuity} of charge $\Div\vJ+\partial_t\rho=0$ as a consequence.
Materials enter only through the three constitutive coefficients
($\vD=\varepsilon\vE$, $\vB=\mu\vH$, $\vJ=\sigma\vE$), scalar for isotropic
media and tensor for anisotropic ones. At interfaces, tangential $\vE$ and
normal $\vB$ are continuous; PEC walls force $\vn\times\vE=\bm0$ (an essential
condition), PMC and symmetry walls impose $\vn\times\vH=\bm0$ for free (a
natural condition). The field stores energy at density
$\tfrac12\varepsilon\abs{\vE}^2+\tfrac1{2\mu}\abs{\vB}^2$, and Poynting's
theorem $\partial_t u+\Div\vS=-\vE\cdot\vJ$ accounts for its flow and loss---the
integrand of the output reductions. The physics divides into three regimes by
the fate of $\partial_t$: dropping it gives the static Poisson and curl--curl
problems posed via the potentials $\vE=-\Grad V$ and $\vB=\Curl\vA$;
substituting $\partial_t\to i\omega$ gives the time-harmonic curl--curl operator
$\mat K+i\omega\mat C-\omega^2\mat M$ behind the eigenmode and driven analyses;
keeping it gives the marched transient system. Every operator the rest of
Part~II assembles is one of these reductions.

\furtherreading
For the physics at a sophomore's level, Griffiths~\citep{griffiths2017} is the
standard and the clearest---its treatment of the four laws, the constitutive
relations, boundary conditions, and Poynting's theorem is exactly the
background assumed here. For the engineering-oriented frequency-domain picture
(phasors, the curl--curl/Helmholtz equation, waveguides) and the finite-element
formulation that \cref{ch:weak} begins, Jin~\citep{jin2014} is the canonical
reference. The mathematically careful treatment of the function spaces
$\Hcurl,\Hdiv$ and the boundary conditions that motivate them---material we
will need in earnest in \cref{ch:functionspaces}---is developed rigorously by
Monk~\citep{monk2003}.

\exercisessection
\begin{enumerate}[nosep]
  \item \textbf{Reading the laws.} For each of the four Maxwell equations, state
  in one sentence what physical quantity is the ``source'' on the right-hand
  side and what field property (divergence or curl) it controls. Which two
  equations couple $\vE$ and $\vB$, and which two are instantaneous
  constraints?

  \item \textbf{Continuity by hand.} Starting from Faraday's law
  $\Curl\vE=-\partial_t\vB$, take the divergence of both sides and use
  $\Div(\Curl\,\cdot)=0$ to show that $\partial_t(\Div\vB)=0$---so once
  $\Div\vB=0$ holds it holds forever. Explain why this is the magnetic analogue
  of the continuity equation~\eqref{eq:continuity}.

  \item \textbf{A different dielectric.} Repeat the capacitance calculation of
  Worked example~\ref{wex:capacitor} for two plates of area $A=\SI{1}{\milli\meter\squared}$
  separated by $d=\SI{10}{\micro\meter}$ of a material with $\varepsilon_r=10$.
  Find $C$, then the energy stored at $V_0=\SI{5}{\volt}$, and verify
  $C=2U/V_0^2$.

  \item \textbf{The displacement current matters.} Suppose Maxwell had written
  Amp\`ere's law without the $\partial_t\vD$ term. Take the divergence and show
  that the resulting equation forces $\Div\vJ=0$ everywhere. Give a physical
  situation (e.g.\ a charging capacitor) where $\Div\vJ\neq0$, and explain why
  the missing term is therefore essential.

  \item \textbf{Refraction, reversed.} Worked example~\ref{wex:refract} sent a
  field from air into a denser dielectric. Reverse it: a field inside the
  dielectric ($\varepsilon_2=4\varepsilon_0$) makes $\SI{66.6}{\degree}$ with the
  normal and exits into air ($\varepsilon_1=\varepsilon_0$). Use the refraction
  law~\eqref{eq:refract} to find the exit angle, and confirm it is the
  $\SI{30}{\degree}$ you started with---refraction is reversible. Then explain,
  in one sentence, why a field line crossing into a \emph{higher}-permittivity
  region always bends away from the normal.

  \item \textbf{The phasor factor.} Show explicitly that if
  $f(t)=\Re\{\hat f\,e^{i\omega t}\}$ then $\partial_t f=\Re\{i\omega\hat
  f\,e^{i\omega t}\}$, justifying the substitution
  $\partial_t\to i\omega$~\eqref{eq:phasor-sub}. What does the factor $i$ do to
  the \emph{phase} of the oscillation, and why does that make ``a derivative''
  into ``a $\SI{90}{\degree}$ phase shift''?

  \item \textbf{From harmonic to static.} Take the harmonic curl--curl equation
  \eqref{eq:curlcurl-harmonic} and let $\omega\to0$. Which term vanishes? Show
  that with a current source restored you recover the magnetostatic equation
  \eqref{eq:curlcurl-static}, confirming that the static regime is the
  zero-frequency limit of the harmonic one.

  \item \textbf{Energy in a resonant mode.} In a lossless cavity a mode
  oscillates with the electric and magnetic energies exchanging back and forth,
  $U_E(t)+U_B(t)=U_{\text{total}}=\text{const}$. Using Poynting's theorem
  \eqref{eq:poynting} with $\vJ=\bm0$ and no net outflow through the PEC walls,
  argue that the total stored energy is conserved. Then add a small
  conductivity $\sigma$ and explain, from the loss term $-\sigma\abs{\vE}^2$,
  why the mode now decays---and why that decay rate is what the quality factor
  $Q$ measures.
\end{enumerate}
